\section{The field comparison and norm detector}\label{sec:field}

Proposition~\ref{prop:integral} constructs an integral class $c$ with
zero generic restriction. To prove it nonzero, we will test its
coefficient image by proper pushforward to $L$ and a valuation.
Localization on the projective curve will introduce arbitrary classes
in $K_3(E;\Z/5)$, for finite extensions $E/L$ arising as
closed-point fields.
We therefore need a description of all these groups that identifies
field transfer with norm.

Keep the class $\beta$ and its restrictions $\beta_E$ fixed as in
Section~\ref{sec:integral}. For every finite extension $E/L$, let
$k_E$ be its residue field, and normalize $v_E$ so that a uniformizer
has valuation one.

\begin{lemma}[Multiplicative field comparison]\label{lem:field}
For $L=\Q_5(\zeta)$ and the fixed $\beta$ above, every finite
extension $E/L$ has an isomorphism
\begin{equation}\label{eq:theta}
 \theta_E:E^\times/E^{\times5}\xrightarrow{\sim}K_3(E;\Z/5),
 \qquad u\bmod E^{\times5}\longmapsto\beta_E\cdot[u].
\end{equation}
The degree-three motivic edge identifies the target with
$H^1_{\mathrm{et}}(E,\mu_5^{\otimes2})$. Under this edge,
the displayed map is Kummer reduction followed by cup product
with the nonzero weight-one edge image
$\alpha_E\in H^0(E,\mu_5)$ of $\beta_E$.
\end{lemma}

\begin{proof}
We first identify the degree-three group, then locate the image of
$\beta_E$ in degree two, and finally compute multiplication by
$\beta_E$ on integral $K_1$. The strongly convergent motivic-to-$K$
spectral sequence for a characteristic-zero field is
\begin{equation}\label{eq:motivic}
 E_2^{a,b}=H^{a-b}(E,\Z/5(-b))
 \quad\Longrightarrow\quad K_{-a-b}(E;\Z/5).
\end{equation}
It has a pairing with the integral spectral sequence that induces
the integral/coefficient $K$-module product on the abutment
\cite[Theorem 16.2 and Appendix C.2--C.3]{FriedlanderSuslin}.
The comparison with motivic complexes preserves products
\cite[Lemma 3.4 and Theorem 6.1]{Suslin}.
Write $s=a-b$ and $j=-b$, so a term $H^s(E,\Z/5(j))$ has total
$K$-degree $2j-s$. Negative weights vanish. For $j\geq0$, the
field motivic groups vanish for $s>j$, and the norm-residue and
motivic--\'etale comparison theorems identify them, for every
$s\leq j$, with
\[
 H^s_{\mathrm{et}}(E,\mu_5^{\otimes j});
\]
see \cite[Theorems 6.16--6.17]{Voevodsky} and
\cite[VI.4.1--4.2]{Weibel}. In particular, the groups vanish for
$s<0$. These comparisons apply because $5$ is invertible in $E$.

The $5$-cohomological dimension of a finite extension of $\Q_5$
is two \cite[VI.\S7, before Proposition 7.3]{Weibel}. Thus the
only potentially nonzero terms satisfy $0\leq s\leq\min(j,2)$.
For $r\geq2$ a differential changes $(s,j)$ to
$(s+2r-1,j+r-1)$, so none can join two nonzero terms. Total
degree three has the single term $(s,j)=(1,2)$, yielding the
natural identification
\begin{equation}\label{eq:degree-three-edge}
 K_3(E;\Z/5)\simeq H^1_{\mathrm{et}}(E,\mu_5^{\otimes2}).
\end{equation}
In particular, this $K$-group is killed by five. In total degree
two the weight-one quotient of the filtration is $H^0(E,\mu_5)$;
the only other possible piece is $H^2(E,\mu_5^{\otimes2})$ in
weight two.

Let $\alpha_E$ be the weight-one image of $\beta_E$.
To see that it is nonzero, restrict to an algebraic closure
$\overline E$. The same comparison there has only $s=0$ terms,
so its total degree-two group has just the weight-one piece
$H^0(\overline E,\mu_5)$. Matsumoto's presentation of field
$K_2$ (\cite[\S5]{Matsumoto}; see also \cite[III.6.1]{Weibel}) gives
$K_2(\overline E)/5=0$: every symbol
$\{u,v\}$ equals $5\{r,v\}$ after choosing $r^5=u$.
The coefficient sequence therefore identifies
$K_2(\overline E;\Z/5)$ with $K_1(\overline E)[5]$ through
$\partial_{\mathrm{coef}}$. The restriction of $\beta_E$ is
nonzero, since its boundary is the primitive root $[\zeta]$.
Naturality of the filtration shows that $\alpha_E$ has nonzero
image in $H^0(\overline E,\mu_5)$. It is consequently a generator
of $H^0(E,\mu_5)$; only this nonvanishing is needed.

It remains to determine the actual product from these filtration
pieces. Integral $K_1(E)=E^\times$ has only the weight-one piece
$H^1(E,\Z(1))$: on the diagonal $2j-s=1$, weights $j\geq2$
give $s>j$ and vanish, weight zero gives
$H^{-1}(E,\Z(0))=0$, and negative weights vanish. The
identification $\Z(1)\simeq\mathbb G_m[-1]$ identifies reduction
of $[u]$ with its Kummer class in
$H^1(E,\mu_5)=E^\times/E^{\times5}$.

Write $F^j$ for the decreasing weight filtration. The preceding
calculations give
\[
 F^1K_1(E)=K_1(E),\quad F^2K_1(E)=0,\qquad
 F^2K_3(E;\Z/5)=K_3(E;\Z/5),\quad F^3K_3(E;\Z/5)=0.
\]
The class $\beta_E$ lies in $F^1K_2(E;\Z/5)$ and has image
$\alpha_E$ in its quotient by $F^2$. The mixed filtered pairing
satisfies
\[
 F^2K_2(E;\Z/5)\cdot F^1K_1(E)
 \subseteq F^3K_3(E;\Z/5)=0.
\]
Thus the possible weight-two part of $\beta_E$ contributes nothing:
the product
$\beta_E\cdot[u]$ is determined by $\alpha_E$ and the
weight-one class of $[u]$. Product compatibility identifies it,
under \eqref{eq:degree-three-edge}, with
$\alpha_E\cup[u]$.
Since $\alpha_E$ generates the constant one-dimensional
$\F_5$-module $\mu_5$, this cup product is an isomorphism.
The product factors through fifth powers because its target is
killed by five and $[u^5]=5[u]$. This proves
\eqref{eq:theta} with its asserted multiplicative identification.
\end{proof}

In particular, the coefficient class already constructed on the
divisor is $a_D=\theta_D(t)$. We now relate the comparison maps for
different extensions of $L$.

\begin{lemma}[Norm detector]\label{lem:norm-detector}
For every finite extension $E/L$
and every $u\in E^\times$,
\begin{equation}\label{eq:transfer}
 \Tr_{E/L}\theta_E(u)=\theta_L(N_{E/L}u).
\end{equation}
Consequently
\[
 \lambda_L=(v_L\bmod5)\circ\theta_L^{-1}:
 K_3(L;\Z/5)\longrightarrow\Z/5
\]
satisfies
\[
 \lambda_L\bigl(\Tr_{E/L}(a)\bigr)=0
 \quad\text{for every }a\in K_3(E;\Z/5)
 \quad\text{if }5\mid[k_E:k_L].
\]
\end{lemma}

\begin{proof}
Restriction of scalars obeys the tensor projection formula at the
level of $K$-theory spectra \cite[V.3.3.2]{Weibel}. Smashing its
base-field factor with $M_5$ gives the mixed pairing identity
\[
 \Tr_{E/L}(\beta_E\cdot[u])
 =\beta\cdot\Tr_{E/L}[u]
 =\beta\cdot[N_{E/L}u].
\]
The second equality uses the norm description of field $K_1$
transfer \cite[III.1.7.1 and the following paragraph]{Weibel}. This proves
\eqref{eq:transfer}.

For normalized valuations, the norm of the prime ideal of $E$
is the $[k_E:k_L]$th power of the prime ideal of $L$. Hence
\[
 v_L(N_{E/L}u)=[k_E:k_L]v_E(u).
\]
Every $a\in K_3(E;\Z/5)$ has the form $\theta_E(u)$ by
Lemma~\ref{lem:field}. Applying the formula to that $u$ proves
the assertion for all $a$. The factor is the residue degree,
not the total degree $[E:L]$.
\end{proof}
